\SourceLeaf{62}motu continuo descripta supponitur, eius elementa $\mathrm{d}P$ (perinde ut elementa superficiei $\mathrm{d}U$) semper positive accipiemus. Cosinus angulorum, quos directio elementi $\mathrm{d}P$ facit cum axibus coordinatarum $x,y,z$, per $X,Y,Z$ denotamus: ne vero sensus directionis ambiguus maneat, hanc ita decernimus, ut ipsa primo loco, directio normalis in elementum $\mathrm{d}P$ superficiem $U$ tangentis et huius respectu introrsum ductae secundo loco, denique normalis in superficiem respectuque spatii $s$ extrorsum ducta tertio loco, constituant systema trium rectarum similiter deinceps sitarum, ut axes coordinatarum $x,y,z$. Ita facile perspicietur (cf. \textit{Disquiss. gen. circa superficies curvas} art. 2), cosinus angulorum inter directionem illam secundam atque axes coordinatarum $x,y,z$ esse resp.
\[
\eta^0Z-\zeta^0Y,\qquad \zeta^0X-\xi^0Z,\qquad \xi^0Y-\eta^0X
\]
si $\xi^0,\eta^0,\zeta^0$ sint valores ipsarum $\xi,\eta,\zeta$ pro puncto elementi $\mathrm{d}P$.

\Article{21}
His ita praeparatis supponamus, superficiem $U$ pati mutationem qualemcunque infinite parvam. Si sufficeret, tales tantummodo mutationes considerare, pro quibus limes $P$ semper invariatus. vel saltem in eadem superficie verticali maneret, manifesto soli coordinatae tertiae $z$ variationem inducere oporteret, quo pacto problema longe facilius evaderet; sed quum problema maxima generalitate nobis ventilandum sit, in tali investigationis modo consideratio variabilitatis limitum in ambages incommodas concinnitatemque turbantes perduceret; quamobrem praestabit, statim ab initio omnes tres coordinatas variationi subiicere. Rem itaque sic imaginabimur, ut cuivis puncto superficiei, cuius coordinatae sunt $x,y,z$, substituamus aliud, cuius coordinatae sint $x+\delta x$, $y+\delta y$, $z+\delta z$, ubi $\delta x,\delta y,\delta z$ spectari possunt tamquam functiones indeterminatae ipsarum $x,y$, sed quarum valores manent infinite parvae. Inquiramus nunc in variationes singulorum elementorum expressionis $W$, et quidem initium faciamus a variatione ipsius elementi $\mathrm{d}U$.

Concipiamus elementum superficiei $U$ triangulare $\mathrm{d}U$ inter puncta, quorum coordinatae sint
\[
\begin{array}{lll}
x,&y,&z\\
x+\mathrm{d}x,&y+\mathrm{d}y.&z+\frac{\mathrm{d}z}{\mathrm{d}x}.\mathrm{d}x+\frac{\mathrm{d}z}{\mathrm{d}y}.\mathrm{d}y\\
x+\mathrm{d}'x,&y+\mathrm{d}'y,&z+\frac{\mathrm{d}z}{\mathrm{d}x}.\mathrm{d}'x+\frac{\mathrm{d}z}{\mathrm{d}y}.\mathrm{d}'y
\end{array}
\]
\SourceLeaf{63}Area duplex huius trianguli per principia nota invenitur
\[
=(\mathrm{d}x.\mathrm{d}'y-\mathrm{d}y.\mathrm{d}'x)\BareRoot{[1+(\frac{\mathrm{d}z}{\mathrm{d}x})^2+(\frac{\mathrm{d}z}{\mathrm{d}y})^2]}
\]
si, quod licet, supponimus, $\mathrm{d}x.\mathrm{d}'y-\mathrm{d}y.\mathrm{d}'x$ esse quantitatem positivam.

In superficie variata loco illorum punctorum tria alia habebimus, quorum coordinatae erunt
\[
\begin{array}{ll}
\text{puncti primi}&x+\delta x,\quad y+\delta y,\quad z+\delta z
\end{array}
\]
puncti secundi
\[
\begin{gathered}
x+\mathrm{d}x+\delta x+\frac{\mathrm{d}\delta x}{\mathrm{d}x}.\mathrm{d}x+\frac{\mathrm{d}\delta x}{\mathrm{d}y}.\mathrm{d}y\\
y+\mathrm{d}y+\delta y+\frac{\mathrm{d}\delta y}{\mathrm{d}x}.\mathrm{d}x+\frac{\mathrm{d}\delta y}{\mathrm{d}y}.\mathrm{d}y\\
z+\frac{\mathrm{d}z}{\mathrm{d}x}.\mathrm{d}x+\frac{\mathrm{d}z}{\mathrm{d}y}.\mathrm{d}y+\delta z+\frac{\mathrm{d}\delta z}{\mathrm{d}x}.\mathrm{d}x+\frac{\mathrm{d}\delta z}{\mathrm{d}y}.\mathrm{d}y
\end{gathered}
\]
puncti tertii
\[
\begin{gathered}
x+\mathrm{d}'x+\delta x+\frac{\mathrm{d}\delta x}{\mathrm{d}x}.\mathrm{d}'x+\frac{\mathrm{d}\delta x}{\mathrm{d}y}.\mathrm{d}'y\\
y+\mathrm{d}'y+\delta y+\frac{\mathrm{d}\delta y}{\mathrm{d}x}.\mathrm{d}'x+\frac{\mathrm{d}\delta y}{\mathrm{d}y}.\mathrm{d}'y\\
z+\frac{\mathrm{d}z}{\mathrm{d}x}.\mathrm{d}'x+\frac{\mathrm{d}z}{\mathrm{d}y}.\mathrm{d}'y+\delta z+\frac{\mathrm{d}\delta z}{\mathrm{d}x}.\mathrm{d}'x+\frac{\mathrm{d}\delta z}{\mathrm{d}y}.\mathrm{d}'y
\end{gathered}
\]
Area duplex trianguli inter haec puncta invenitur per eandem methodum
\[
=(\mathrm{d}x.\mathrm{d}'y-\mathrm{d}y.\mathrm{d}'x)\BareRoot{N}
\]
si brevitatis caussa per $N$ denotatur aggregatum
\[
\begin{aligned}
&[(1+\frac{\mathrm{d}\delta x}{\mathrm{d}x})(1+\frac{\mathrm{d}\delta y}{\mathrm{d}y})-\frac{\mathrm{d}\delta x}{\mathrm{d}y}.\frac{\mathrm{d}\delta y}{\mathrm{d}x}]^2\\
+&[(1+\frac{\mathrm{d}\delta x}{\mathrm{d}x})(\frac{\mathrm{d}z}{\mathrm{d}y}+\frac{\mathrm{d}\delta z}{\mathrm{d}y})-\frac{\mathrm{d}\delta x}{\mathrm{d}y}(\frac{\mathrm{d}z}{\mathrm{d}x}+\frac{\mathrm{d}\delta z}{\mathrm{d}x})]^2\\
+&[(1+\frac{\mathrm{d}\delta y}{\mathrm{d}y})(\frac{\mathrm{d}z}{\mathrm{d}x}+\frac{\mathrm{d}\delta z}{\mathrm{d}x})-\frac{\mathrm{d}\delta y}{\mathrm{d}x}(\frac{\mathrm{d}x}{\mathrm{d}y}+\frac{\mathrm{d}\delta z}{\mathrm{d}y})]^2
\end{aligned}
\]
Facta evolutione et reiectis quantitatibus secundi ordinis, invenitur
\[
\BareRoot{N}=\BareRoot{[1+(\frac{\mathrm{d}z}{\mathrm{d}x})^2+(\frac{\mathrm{d}z}{\mathrm{d}y})^2]}.\left[1+\frac{L}{1+(\frac{\mathrm{d}z}{\mathrm{d}x})^2+(\frac{\mathrm{d}z}{\mathrm{d}y})^2}\right]
\]
si brevitatis gratia per $L$ denotatur aggregatum
\SourceLeaf{64}\[
\begin{aligned}
&\frac{\mathrm{d}\delta x}{\mathrm{d}x}[1+(\frac{\mathrm{d}z}{\mathrm{d}y})^2]-\frac{\mathrm{d}\delta x}{\mathrm{d}y}.\frac{\mathrm{d}z}{\mathrm{d}x}.\frac{\mathrm{d}z}{\mathrm{d}y}-\frac{\mathrm{d}\delta y}{\mathrm{d}x}.\frac{\mathrm{d}z}{\mathrm{d}x}.\frac{\mathrm{d}z}{\mathrm{d}y}\\
&+\frac{\mathrm{d}\delta y}{\mathrm{d}y}[1+(\frac{\mathrm{d}z}{\mathrm{d}x})^2]+\frac{\mathrm{d}\delta z}{\mathrm{d}x}.\frac{\mathrm{d}z}{\mathrm{d}x}+\frac{\mathrm{d}\delta z}{\mathrm{d}y}.\frac{\mathrm{d}z}{\mathrm{d}y}
\end{aligned}
\]
Est itaque ratio trianguli primi ad secundum ut $1$ ad
\[
1+\frac{L}{1+(\frac{\mathrm{d}z}{\mathrm{d}x})^2+(\frac{\mathrm{d}z}{\mathrm{d}y})^2}
\]
adeoque independens a figura trianguli $\mathrm{d}U$, resultatque
\[
\delta\mathrm{d}U=\frac{L\mathrm{d}U}{1+(\frac{\mathrm{d}z}{\mathrm{d}x})^2+(\frac{\mathrm{d}z}{\mathrm{d}y})^2}
\]
sive in terminis explicitis
\[
\begin{aligned}
\delta\mathrm{d}U=\mathrm{d}U\big(&\frac{\mathrm{d}\delta x}{\mathrm{d}x}(\eta\eta+\zeta\zeta)-\frac{\mathrm{d}\delta x}{\mathrm{d}y}.\xi\eta-\frac{\mathrm{d}\delta y}{\mathrm{d}x}.\xi\eta\\
&+\frac{\mathrm{d}\delta y}{\mathrm{d}y}(\xi\xi+\zeta\zeta)-\frac{\mathrm{d}\delta z}{\mathrm{d}x}.\xi\zeta-\frac{\mathrm{d}\delta z}{\mathrm{d}y}.\eta\zeta\big)
\end{aligned}
\]

\Article{22}
Variationem totius superficiei $U$ obtinebimus per integrationem huius expressionis per omnia elementa $\mathrm{d}U$ extendendam. Ad hunc finem duas huius integralis partes, puta
\[
\int\mathrm{d}U.((\eta\eta+\zeta\zeta)\frac{\mathrm{d}\delta x}{\mathrm{d}x}-\xi\eta\frac{\mathrm{d}\delta y}{\mathrm{d}x}-\xi\zeta\frac{\mathrm{d}\delta z}{\mathrm{d}x})=A
\]
atque
\[
\int\mathrm{d}U(-\xi\eta\frac{\mathrm{d}\delta x}{\mathrm{d}y}+(\xi\xi+\zeta\zeta)\frac{\mathrm{d}\delta y}{\mathrm{d}y}-\eta\zeta\frac{\mathrm{d}\delta z}{\mathrm{d}y})=B
\]
seorsim tractabimus.

Concipiatur planum axi coordinatarum $y$ normale, et quidem tale, ut valor determinatus ipsius $y$, ei competens, sit intra ambitum valorum extremorum, quos habet $y$ in superficie $U$. Hoc planum peripheriam $P$ secabit vel in duobus, vel in quatuor, vel in sex etc. punctis, quorum coordinatae primae sint deinceps $x^0,x',x''$ etc.; perinde reliquae quantitates ad haec puncta pertinentes per indices distinguantur. Eodem modo secetur superficies per aliud planum illi infinite propinquum et parallelum, cui competat coordinata secunda $y+\mathrm{d}y$; inter haec plana reperientur elementa peripheriae $\mathrm{d}P^0$, $\mathrm{d}P'$, $\mathrm{d}P''$ etc., perspicieturque facile, haberi
\[
\mathrm{d}y=-Y^0\mathrm{d}P^0=+Y'\mathrm{d}P'=-Y''\mathrm{d}P''=+Y'''\mathrm{d}P'''\quad\text{etc.}
\]
\SourceLeaf{65}Si insuper concipimus infinite multa plana axi coordinatarum $x$ normalia, cuivis elemento $\mathrm{d}x$ inter $x^0$ et $x'$, vel inter $x''$ et $x'''$ etc. sito respondebit elementum $\mathrm{d}U=\frac{\mathrm{d}x.\mathrm{d}y}{\zeta}$, unde patet, eam partem integralis $A$, quae respondet parti superficiei inter plana $y,y+\mathrm{d}y$ sitae, haberi ex integratione
\[
\mathrm{d}y\int\mathrm{d}x.(\frac{\eta\eta+\zeta\zeta}{\zeta}.\frac{\mathrm{d}\delta x}{\mathrm{d}x}-\frac{\xi\eta}{\zeta}.\frac{\mathrm{d}\delta y}{\mathrm{d}x}-\xi\frac{\mathrm{d}\delta z}{\mathrm{d}x})
\]
extensa ab $x=x^0$ usque ad $x=x'$, dein ab $x=x''$ usque ad $x=x'''$ etc. Indefinite vero hoc integrale exhibetur per
\[
\begin{aligned}
&(\frac{\eta\eta+\zeta\zeta}{\zeta}.\delta x-\frac{\xi\eta}{\zeta}.\delta y-\xi\delta z)\mathrm{d}y\\
&-\mathrm{d}y\int(\delta x.\frac{\mathrm{d}\frac{\eta\eta+\zeta\zeta}{\zeta}}{\mathrm{d}x}-\delta y.\frac{\mathrm{d}\frac{\xi\eta}{\zeta}}{\mathrm{d}x}-\delta z.\frac{\mathrm{d}\xi}{\mathrm{d}x})\mathrm{d}x
\end{aligned}
\]
unde colligitur, prodire pro casu nostro
\[
\begin{aligned}
&(\frac{\eta^0\eta^0+\zeta^0\zeta^0}{\zeta^0}.\delta x^0-\frac{\xi^0\eta^0}{\zeta^0}.\delta y^0-\xi^0\delta z^0)Y^0\mathrm{d}P^0\\
+&(\frac{\eta'\eta'+\zeta'\zeta'}{\zeta'}.\delta x'-\frac{\xi'\eta'}{\zeta'}.\delta y'-\xi'\delta z')Y'\mathrm{d}P'\\
+&(\frac{\eta''\eta''+\zeta''\zeta''}{\zeta''}.\delta x''-\frac{\xi''\eta''}{\zeta''}.\delta y''-\xi''\delta z'')Y''\mathrm{d}P''\\
+&\text{etc.}\\
-&\int\zeta\mathrm{d}U.(\delta x.\frac{\mathrm{d}\frac{\eta\eta+\zeta\zeta}{\zeta}}{\mathrm{d}x}-\delta y.\frac{\mathrm{d}\frac{\xi\eta}{\zeta}}{\mathrm{d}x}-\delta z.\frac{\mathrm{d}\xi}{\mathrm{d}x})
\end{aligned}
\]
sive, quod idem est,
\[
\begin{aligned}
&\Sigma(\frac{\eta\eta+\zeta\zeta}{\zeta}.\delta x-\frac{\xi\eta}{\zeta}.\delta y-\xi\delta z)Y\mathrm{d}P\\
&-\int\zeta\mathrm{d}U.(\delta x.\frac{\mathrm{d}\frac{\eta\eta+\zeta\zeta}{\zeta}}{\mathrm{d}x}-\delta y.\frac{\mathrm{d}\frac{\xi\eta}{\zeta}}{\mathrm{d}x}-\delta z.\frac{\mathrm{d}\xi}{\mathrm{d}x})
\end{aligned}
\]
ubi tum summatio per omnia elementa $\mathrm{d}P$, tum integratio per omnia elementa $\mathrm{d}U$, intra plana $y$ et $y+\mathrm{d}y$ sita, extendenda est.

Tota itaque quantitas $A$ exprimetur per
\[
\begin{aligned}
&\int(\frac{\eta\eta+\zeta\zeta}{\zeta}.\delta x-\frac{\xi\eta}{\zeta}.\delta y-\xi\delta z)Y\mathrm{d}P\\
&-\int\zeta\mathrm{d}U(\delta x.\frac{\mathrm{d}\frac{\eta\eta+\zeta\zeta}{\zeta}}{\mathrm{d}x}-\delta y.\frac{\mathrm{d}\frac{\xi\eta}{\zeta}}{\mathrm{d}x}-\delta z.\frac{\mathrm{d}\xi}{\mathrm{d}x})
\end{aligned}
\]
ubi integrationem priorem per totam peripheriam $P$, posteriorem per totam superficiem $U$ extendere oportet.

\SourceLeaf{66}\Article{23}
Per ratiocinia prorsus similia invenimus
\[
\begin{aligned}
B={}&\int(\frac{\xi\eta}{\zeta}.\delta x-\frac{\xi\xi+\zeta\zeta}{\zeta}.\delta y+\eta\delta z)X\mathrm{d}P\\
&+\int\zeta\mathrm{d}U(\delta x.\frac{\mathrm{d}\frac{\xi\eta}{\zeta}}{\mathrm{d}y}-\delta y.\frac{\mathrm{d}\frac{\xi\xi+\zeta\zeta}{\zeta}}{\mathrm{d}y}+\delta z.\frac{\mathrm{d}\eta}{\mathrm{d}y})
\end{aligned}
\]
Statuendo itaque, pro quovis puncto peripheriae $P$,
\[
\begin{aligned}
&[X\xi\eta+Y(\eta\eta+\zeta\zeta)]\delta x-[X(\xi\xi+\zeta\zeta)+Y\xi\eta]\delta y\\
&+(X\eta\zeta-Y\xi\zeta)\delta z=\zeta Q
\end{aligned}
\]
nec non, pro quovis puncto superficiei $U$,
\[
\begin{aligned}
&(\frac{\mathrm{d}\frac{\xi\eta}{\zeta}}{\mathrm{d}y}-\frac{\mathrm{d}\frac{\eta\eta+\zeta\zeta}{\zeta}}{\mathrm{d}x})\zeta\delta x+(\frac{\mathrm{d}\frac{\xi\eta}{\zeta}}{\mathrm{d}x}-\frac{\mathrm{d}\frac{\xi\xi+\zeta\zeta}{\zeta}}{\mathrm{d}y})\zeta\delta y\\
&+(\frac{\mathrm{d}\xi}{\mathrm{d}x}+\frac{\mathrm{d}\eta}{\mathrm{d}y})\zeta\delta z=V
\end{aligned}
\]
erit tandem
\[
\delta U=\int Q\mathrm{d}P+\int V\mathrm{d}U
\]
ubi integratio prima per totam peripheriam $P$, secunda per totam superficiem $U$ extendi debet.

\Article{24}
Formulas pro $Q$ et $V$ modo allatas notabiliter contrahere licet. Et quidem, adiumento aequationis $X\xi+Y\eta+Z\zeta=0$, $Q$ statim induit formam symmetricam sequentem:
\[
Q=(Y\zeta-Z\eta)\delta x+(Z\xi-X\zeta)\delta y+(X\eta-Y\xi)\delta z
\]

Quo etiam expressio pro $V$ eruta in formam concinniorem reducatur, observamus, e formulis
\[
\frac{\mathrm{d}z}{\mathrm{d}x}=-\frac{\xi}{\zeta},\qquad \frac{\mathrm{d}z}{\mathrm{d}y}=-\frac{\eta}{\zeta}
\]
sequi
\[
\frac{\mathrm{d}\frac{\xi}{\zeta}}{\mathrm{d}y}=\frac{\mathrm{d}\frac{\eta}{\zeta}}{\mathrm{d}x}
\]
Hinc fit
\[
\frac{\mathrm{d}\frac{\xi\eta}{\zeta}}{\mathrm{d}y}=\frac{\xi}{\zeta}.\frac{\mathrm{d}\eta}{\mathrm{d}y}+\eta\frac{\mathrm{d}\frac{\xi}{\zeta}}{\mathrm{d}y}=\frac{\xi}{\zeta}.\frac{\mathrm{d}\eta}{\mathrm{d}y}+\eta\frac{\mathrm{d}\frac{\eta}{\zeta}}{\mathrm{d}x}
\]
